\chapter{Do Inferred Specifications Help LLMs Write Better Tests?}\label{ch:article1}

This chapter establishes the foundational premise on which the remainder of the
thesis rests: that automatically inferred specifications, despite being heuristic
and partial, are good enough to be useful in a concrete downstream task. It asks
whether supplying inferred contracts to a large language model measurably improves
the unit tests the model writes, and answers in the affirmative through a
controlled comparison on a mature Java library. In doing so it converts the
abstract value of specifications --- developed conceptually in \cref{ch:background}
--- into measured, downstream utility, and so supplies the evidence for RQ2 of
\cref{ch:introduction}.

\section{Motivation}
\label{sec:art1:motivation}

Large language models have rapidly become a practical means of generating unit
tests from source code~\citep{chen2021codex,tufano2020athenatest,yuan2024chattester,schafer2024testpilot}.
Given a method signature, a modern model produces compilable, plausibly named
tests in seconds; given the source body, it produces tests that exercise the
visible control-flow paths. The persistent weakness is the \emph{oracle problem}~\citep{barr2015oracle}:
the generated tests assert what the implementation does, not what it should do.
When the model is given only a signature it has no oracle at all and falls back on
weak assertions such as \enquote{does not throw}; when it is given the body it
tends to encode the program's actual behaviour rather than its intended
behaviour, so that a bug in the implementation is dutifully recorded as the
expected result. Recent oracle-generation work confirms this tendency
empirically~\citep{liu2026oracle,watson2020atlas}, and large-scale studies of
LLM-generated code report that fluency does not entail correctness: hallucinated
or subtly wrong outputs are common even when the surface form is
plausible~\citep{zhang2025hallucination,liu2023evalplus}.

The difficulty is structural rather than incidental. A unit test has two parts: a
prefix that drives the program into some state, and an oracle that asserts what
should then be true. A language model is good at the prefix --- constructing
arguments, instantiating objects, arranging calls --- because that is a matter of
syntactic fluency, and fluency is what these models have in abundance. The oracle
is the hard part, because it requires knowing what the method is \emph{supposed} to
do, and that knowledge is not present in a signature and is only implicitly, and
unreliably, present in a body. Faced with a signature alone, the model has nothing
to assert and retreats to assertions that hold of almost any implementation; faced
with a body, it has too much, and asserts whatever the body happens to compute,
bugs included. Either way the oracle the model produces is weak or wrong for the
same underlying reason: it has no independent statement of intent to assert
against. A specification is precisely such a statement.

Formal specifications are the classical answer to the oracle
problem~\citep{meyer1992dbc,leavens2006jml}, but their adoption is limited by the
cost of writing them. The empirical question this chapter addresses is therefore
practical rather than theoretical: \emph{does giving a specification to a large
language model measurably improve the unit tests it writes?} The hypothesis is
that a specification --- even one that is heuristically derived and incomplete ---
carries oracle information that is absent from the signature and largely redundant
with the source body, and that this information is dense enough to shift
LLM-generated tests from coverage-driven artefacts towards genuine behavioural
checks.

The notation, tooling, and inference instrument used throughout are those
introduced in \cref{ch:background}: the Java Modeling Language (JML) for the
contract notation; the weakest-precondition (WP) and strongest-postcondition (SP)
correspondences that organise inference; OpenJML's Extended Static Checker (ESC)
for formal validation~\citep{cok2014openjml}; the Z3 SMT solver~\citep{demoura2008z3}
as the underlying decision procedure; and the heuristic inference engine
(hereafter the inferrer) whose construction and validation \cref{ch:background}
describes. This chapter does not re-derive any of that material; it takes the
inferrer as a given supply of specifications and measures their downstream effect.
What follows is therefore deliberately confined to the study itself: the
experimental manipulation, its design, the metrics, the results, and the threats
specific to this experiment.

\subsection{Why test generation is the measurable case}
\label{sec:art1:why-tests}

Inferred specifications have several downstream uses beyond LLM test generation,
and the thesis pursues them in later chapters: they act as the callee summaries
that compositional verification of calling code requires (\cref{ch:article3}); they
serve as a machine-readable contract that an AI coding agent can consume when
reasoning about a library that postdates its training cut-off (\cref{ch:article5});
and they strengthen the oracles of any specification-aware test-generation tool
that consumes JML, such as EvoSuite~\citep{fraser2011evosuite} or
Randoop~\citep{pacheco2007randoop}, rather than only LLM-driven ones. Test
generation is singled out here not because it is the most important use but
because it is the one whose effect size can be measured cleanly within the budget
of a single study. Mutation testing supplies an objective, automatable proxy for
oracle strength: a test suite that kills more mutants is one whose assertions
distinguish more behavioural deviations, regardless of how the assertions were
authored. The remaining uses are pursued qualitatively or in separate studies; the
measurement here is what licenses the claim, foundational to the thesis, that
inferred specifications are useful and not merely plausible.

\section{Approach}
\label{sec:art1:approach}

The study holds the language model fixed and varies only the material supplied to
it alongside an otherwise identical method signature. Four conditions are compared
on the same corpus of methods, designed to isolate the contribution of formal
specifications to test-generation quality:

\begin{description}
  \item[Phase~1 (P1) --- Signature only.] Test generation from the method
    signature alone. This is the floor: the model has a name and a type, but no
    statement of intended behaviour and no body.
  \item[Phase~2 (P2) --- Signature with guidance.] The signature supplemented with
    a fixed, method-agnostic checklist of edge-case categories the model is
    instructed to cover. This is the control for the alternative hypothesis that
    any explicit guidance, rather than specifications in particular, is what raises
    test quality.
  \item[Phase~3 (P3) --- Signature with specification.] The signature augmented
    with the formal specification inferred by the inferrer for that method, used
    without modification. This is the treatment condition.
  \item[Phase~4 (P4) --- Source-code baseline.] The full method source body,
    providing an upper bound on the quality attainable when the model can see the
    implementation directly.
\end{description}

The four conditions form an ordered ladder of information density. P1 supplies the
least, P4 the most, and P3 supplies a compact behavioural summary that occupies an
order of magnitude fewer lines than the body it summarises. The design is
deliberately parsimonious: the signature is held constant across all four phases,
so any difference in the generated tests is attributable to the material added
alongside it, not to a change in what the model is being asked to test.

\subsection{The role of the guidance control}
\label{sec:art1:p2-control}

Phase~2 exists to defeat a specific rival explanation. One might object that the
benefit attributed to specifications in P3 is really the benefit of any structured
prompt: a model told explicitly to think about edge cases will write more
thorough tests regardless of whether the guidance is method-specific. P2
operationalises that objection. Its guidance is a uniform block appended to the P1
prompt, identical for every method, conveying no per-method behavioural
information. It enumerates the standard edge-case categories of unit-test design:
\texttt{null} for each reference parameter; empty, single-element, and
maximum-size collections, arrays, and strings; boundary numeric values (zero,
$\pm 1$, the integer extrema, and the adjacent integers); floating-point special
values (\texttt{NaN}, $\pm$infinity, $\pm 0.0$); unicode boundary code points for
character and string parameters; and overflow-prone arithmetic combinations. If a
generic edge-case checklist were to match P3 on oracle strength, the contribution
of formal specifications would be indistinguishable from that of any other
structured prompt, and the thesis premise would be unsupported. The result, given
in \cref{sec:art1:results}, is that it does not: P2 generates the most tests of any
phase yet attains a lower mutation score than P3, which is the cleanest single
piece of evidence that test quantity is a poor proxy for oracle strength.

\subsection{A worked example of the manipulation}
\label{sec:art1:worked-example}

To make the experimental manipulation concrete, consider a single method,
\texttt{MutableInt.add(Number)} from the \texttt{org.apache.commons.lang3.mutable}
package. \cref{lst:art1:phase-example} sets out the four phase inputs side by
side. The signature is identical in every phase; only the material supplied
alongside it changes. P2 adds the method-agnostic edge-case checklist; P3 adds the
JML contract emitted by the inferrer for this method, without modification; and P4
substitutes the full source body. The P3 contract states that the operand must be
non-null, that the method assigns to the \texttt{value} field, and that the new
value equals the old value plus the operand's integer value --- a complete oracle
for the method's behaviour, expressed in three lines.

\begin{lstlisting}[language=Java,caption={Worked example of the four phase inputs supplied to the model, applied to \texttt{MutableInt.add(Number)}. The signature is held constant across phases; P2 adds a method-agnostic edge-case checklist, P3 adds the inferred JML contract, and P4 substitutes the full source body.},label={lst:art1:phase-example}]
// P1 -- Signature only:
public void add(Number operand);

// P2 -- Signature + edge-case checklist (identical block for every method):
public void add(Number operand);
// Test-design checklist:
//  - null for each reference parameter
//  - empty/single-element/maximum-size collections, arrays, strings
//  - boundary numeric values (0, +/-1, integer extrema, adjacent integers)
//  - floating-point specials (NaN, +/-Infinity, +/-0.0)
//  - unicode boundary code points (for char/String parameters)
//  - overflow-prone arithmetic combinations

// P3 -- Signature + inferred JML specification:
//@ requires operand != null;
//@ assignable this.value;
//@ ensures this.value == \old(this.value) + operand.intValue();
public void add(Number operand);

// P4 -- Full source code (upper-bound baseline):
public void add(Number operand) {
    this.value += operand.intValue();
}
\end{lstlisting}

\subsection{The specifications the model is given}
\label{sec:art1:specs-given}

The Phase~3 contracts are exactly the output of the inference instrument described
in \cref{ch:background}, used without manual editing; this subsection records only
what is needed to interpret the experiment, namely the shape and provenance of the
material the model sees. Each contract is a small block of JML comprising up to
four kinds of clause: \texttt{requires} clauses (preconditions on the entry state),
\texttt{ensures} clauses (postconditions relating the result and final state to
the entry state), \texttt{assignable} clauses (the frame, naming what the method
may modify), and, for looping methods, \texttt{loop\_invariant} clauses. The
inferrer derives these by pattern matching over the abstract syntax tree, organised
by the weakest-precondition and strongest-postcondition correspondences set out in
\cref{ch:background}: guard-and-throw idioms become preconditions, return
statements and field assignments become postconditions (with field references in
the pre-state wrapped so that the clause refers to the value on entry), and a
family of loop heuristics supplies invariants for counter and accumulator loops.
Branching is preserved at clause granularity: a method that returns one value on
one path and another on a second path yields two guarded postconditions, each
pairing its return value with the condition that reaches it, rather than a single
collapsed disjunction. The practical consequence for the experiment is that the
contract the model receives is decomposable --- the model can pick off one clause
at a time --- which is one of the mechanisms by which Phase~3 attains high
mutant-kill density per test (\cref{sec:art1:why-help}).

Two properties of this supply are what make the experiment well-posed. First, the
contracts are deterministic: the same method always yields the same contract, so
Phase~3 is a controlled input rather than a sampled one. Second, the contracts are
formally validated for consistency with the implementation, so that the
test-generation effect cannot be dismissed as the consequence of feeding the model
inconsistent oracle information; the validation results are reported in
\cref{sec:art1:spec-quality}.

\section{Experimental Design}
\label{sec:art1:design}

\subsection{Subject selection}
\label{sec:art1:subjects}

Apache Commons Lang~3.14 was selected as the evaluation subject for its diversity
of method types, its extensive Javadoc documentation (which supports independent
validation), its maturity, its preponderance of deterministic pure functions
(well suited to mutation testing), and its precedent in the software-engineering
research literature, where it recurs as a benchmark for test-generation
tools~\citep{pacheco2007randoop,fraser2011evosuite}. Eleven classes were selected
to cover a range of method behaviours: \texttt{NumberUtils}, \texttt{BooleanUtils},
and \texttt{CharUtils} (utility helpers); \texttt{MutableInt},
\texttt{MutableDouble}, and \texttt{MutableBoolean} (state modification);
\texttt{Pair} and \texttt{MutablePair} (factory and delegation); \texttt{Validate}
and \texttt{Range} (control flow); and the \texttt{Mutable} interface, which
declares the abstract operations realised by the three \texttt{Mutable*} classes.
All public, package-private, and protected methods of these classes were included;
no method was excluded a~posteriori. The resulting corpus comprises 312 methods.

Each method was assigned to one of six structural categories adapted from a
standard method-stereotype taxonomy: \emph{accessors and mutators},
\emph{control flow}, \emph{factory and delegate}, \emph{utility},
\emph{state modification}, and \emph{other} (event handlers and callbacks). The
distribution across the corpus is given in \cref{tab:art1:method-distribution}.
The \emph{other} category contains zero methods, because Apache Commons Lang is a
library of pure utility and value types and includes no event-handler or callback
methods; the consequence of its absence for external validity is taken up in
\cref{sec:art1:threats}.

\begin{table}[ht]
\centering
\caption{Distribution of methods by structural category in the Apache Commons
Lang subset. The \emph{other} category contains zero methods because the library
includes no event-handler or callback methods.}
\label{tab:art1:method-distribution}
\begin{tabular}{lrr}
\toprule
\textbf{Category} & \textbf{Count} & \textbf{\%} \\
\midrule
Utility methods               & 89  & 28.5\% \\
State modification            & 78  & 25.0\% \\
Accessors \& mutators         & 62  & 19.9\% \\
Control flow                  & 48  & 15.4\% \\
Factory \& delegate           & 35  & 11.2\% \\
Other (callbacks, handlers)   & 0   & 0.0\%  \\
\midrule
\textbf{Total}                & \textbf{312} & \textbf{100\%} \\
\bottomrule
\end{tabular}
\end{table}

\subsection{Model and generation protocol}
\label{sec:art1:protocol}

Test generation used Google's Gemini~2.0, configured with a sampling temperature
of 0.3 to balance diversity against consistency, a top-$p$ of 0.95, and a
maximum of 4{,}096 output tokens. Five independent runs were performed per method
per phase to control for sampling non-determinism, giving $312 \times 4 \times 5 =
6{,}240$ generations for the main experiment; with pilot tuning the total number
of model invocations across the study was approximately 25{,}000. The prompts were
kept minimal and differed across phases only in the inclusion of the edge-case
checklist (P2), the inferred specification (P3), or the source body (P4); the
verbatim templates accompany the replication package.

A single model family was used, for budget reasons and because Gemini~2.0 is
widely available and representative of current instruction-tuned models. The claim
the study defends is the qualitative ordering of the four conditions, not the
absolute scores under any particular model; the corresponding single-model threat
is discussed in \cref{sec:art1:threats}.

The generated tests were evaluated by a fixed, fully automated pipeline so that no
human judgement entered the metrics. For each method and phase, the five generated
test classes were compiled against the unmodified library, the compiling classes
were executed against the reference implementation to obtain a pass rate, and the
passing suite was then run under PIT to obtain a mutation score. Compilation
failures and runtime failures were recorded rather than discarded, since they are
themselves an outcome of interest: a phase that produces many tests but few that
compile is supplying the model with material it cannot turn into valid code, and
the compilation and pass rates in \cref{tab:art1:test-results} capture exactly
that. Equivalent mutants, which inflate the mutation denominator by being
behaviourally indistinguishable from the original, were filtered by PIT's default
heuristics; the residual was estimated by sampling and is small enough not to
disturb the between-phase comparison, which is in any case a difference of
differences in which a constant equivalent-mutant rate cancels.

\subsection{Metrics}
\label{sec:art1:metrics}

Four test-quality metrics are reported. \emph{Test count} is the per-method
average number of generated tests across the five runs. \emph{Compilation rate} is
the proportion of generated test classes that compile. \emph{Pass rate} is the
proportion of compiling tests that pass against the reference implementation.
\emph{Mutation score} is the proportion of injected faults that the generated
suite detects, computed with PIT (version~1.15.3) under its \texttt{DEFAULTS}
operator group, with a per-test timeout of ten seconds. The operator group covers
conditional-boundary mutations (for example replacing \texttt{<} with
\texttt{<=}), increment replacements, negation of arithmetic and boolean
expressions, binary-arithmetic operator swaps, and a family of return-value
substitutions (empty, false, true, null, and primitive returns). Mutation scores
are reported on the post-filter mutant set: PIT's default heuristics remove the
most trivially equivalent mutants, and the residual rate of trivially equivalent
mutants was estimated at 4.7\% by a dual-rater inspection of a random sample of
survived mutants (Cohen's $\kappa = 0.83$).

Mutation score is the headline metric because it is the available objective proxy
for oracle strength~\citep{barr2015oracle}: a suite that kills more mutants is one
whose assertions distinguish more behavioural deviations. Test count is reported
alongside it precisely so that the two can be contrasted; the central finding of
the study turns on a case where the two diverge. The choice of mutation score over
simpler measures is deliberate. Structural coverage --- the proportion of
statements or branches a suite exercises --- is insensitive to assertion strength:
a test that calls a method and asserts nothing about the result still covers the
code it runs, so a suite can reach full coverage while detecting almost no faults.
Mutation testing closes this gap by requiring the suite to \emph{distinguish} a
deliberately altered program from the original, which can only happen if the
assertions encode something about expected behaviour. It is for this reason the
standard yardstick in the test-generation literature, and it is the property the
study most wishes to measure, because oracle strength is exactly what a
specification is hypothesised to supply.

It is worth being explicit about what discharge analysis (\cref{sec:art1:spec-quality})
and the test-generation metrics each measure, because the same 312 methods feed
both. Discharge is a check on the supply of specifications --- whether each
inferred clause is consistent with the implementation it describes --- and is
computed by OpenJML without any reference to the language model. The
test-generation metrics are a check on the downstream effect of that supply ---
whether a model handed the specification writes better tests --- and are computed
without any reference to the discharge outcome. The model is never told which
clauses discharged, and the discharge verdict never depends on what tests the
model wrote. The two are therefore independent observations on one population, and
neither can contaminate the other.

\subsection{Statistical analysis}
\label{sec:art1:stats}

For each metric the analysis reports the mean, standard deviation, and 95\%
bootstrap confidence intervals (10{,}000 iterations). Phase comparisons use paired
$t$-tests, corroborated by Wilcoxon signed-rank tests where the within-phase
distribution departs from normality, with Cohen's $d$ as the effect-size measure.
Because there are four phases, there are $\binom{4}{2} = 6$ pairwise contrasts per
metric and $6 \times 4 = 24$ contrasts in total; a Bonferroni correction is
applied across this family at a family-wise $\alpha = 0.05$, giving a per-test
threshold of $\alpha' \approx 0.0021$. Reported significant differences survive
this corrected threshold and are stated as $p < 0.001$ throughout. A power
analysis based on a 30-method pilot (within-pair standard deviation of 11.4
percentage points) puts the minimum effect detectable at the corrected threshold
with 80\% power at 2.2 percentage points; the observed principal effect is more
than an order of magnitude larger, giving a post-hoc power exceeding~0.999.

\section{Results}
\label{sec:art1:results}

The results fall into two parts. The first characterises the specifications
themselves --- how many methods receive a non-trivial contract, how strong those
contracts are, and what proportion of the inferred clauses OpenJML's Extended
Static Checker can discharge against the implementation. These are properties of
the supply of specifications that Phase~3 consumes, and they bound what the
downstream experiment can possibly demonstrate: a specification that is vacuous or
inconsistent cannot help the model. The second part is the test-generation
experiment proper. The two analyses are methodologically independent observations
on the same corpus. The discharge analysis is a verifier consistency check ---
is the specification formally consistent with the implementation? --- whereas the
effectiveness analysis compares model-generated tests across input conditions ---
does the specification, whatever its accuracy, change what the model produces? No
leakage between the two is possible: the model never sees the discharge outcome,
and the discharge outcome is computed independently of the test-generation
outcome.

\subsection{Coverage, strength, and discharge of the inferred specifications}
\label{sec:art1:spec-quality}

The inferrer emits a non-trivial specification --- at least one non-trivial
precondition or postcondition --- for 309 of the 312 methods, a coverage rate of
99.0\%. The three uncovered methods are overloads of
\texttt{NumberUtils.toScaledBigDecimal} that take a rounding-mode argument and
delegate to \texttt{BigDecimal.setScale}, for which no standard-library
specification entry was available; the engine emits an empty contract rather than
guess, which is the conservative limit of interprocedural summarisation rather
than a heuristic failure. For context, a documentation-based inference baseline
(Jdoctor, which recovers exception preconditions from Javadoc by natural-language
processing) covers 45.2\% of the same corpus, and a direct LLM baseline (the same
Gemini~2.0 model prompted to emit JML from source) covers all 312 methods but with
only 72.4\% run-to-run consistency: prompted identically on the same method, the
model returns non-identical contracts across runs, and the variance includes
clauses a verifier would reject. The static inferrer is deterministic by
construction and so attains 100\% consistency; this determinism is one of the
reasons the study can treat the Phase~3 input as a controlled variable at all.

Specifications are classified by structural strength: \emph{strong} when both a
non-trivial precondition and a non-trivial postcondition are present,
\emph{partial} when exactly one side is non-trivial, and \emph{weak} when neither
is. The classification is computed directly from the emitted JML, independently of
whether the clauses discharge. Strength varies systematically with category.
Accessors and factory methods reach the highest proportion of strong
specifications (78.2\% and 72.1\% respectively), because their behaviour is largely
captured by a single return relationship and a simple guard. State-modification
methods are weakest (58.7\% strong), reflecting the difficulty of capturing a
complete frame condition for a method that mutates object state under conditional
logic. This ordering matters for the downstream result, because it predicts ---
and \cref{sec:art1:by-category} confirms --- that the categories with the strongest
specifications are not always the ones that gain most: an accessor with a strong
specification still gains little downstream because its signature-only baseline is
already high, whereas a control-flow method with a merely partial specification
gains a great deal because its baseline is low.

The clause-level discharge results are summarised in
\cref{tab:art1:discharge}. Across all 985 inferred clauses --- preconditions,
postconditions, loop invariants, and assignable (frame) clauses --- OpenJML's
Extended Static Checker discharges 92.9\% against the implementation. Two
interpretive caveats apply. First, discharge is a \emph{soundness} measure, not a
\emph{completeness} one: a vacuous contract would also discharge, so the figure
should be read as \enquote{no clause was found inconsistent with the
implementation}, which is exactly why the structural strength measure above is
reported alongside it. Second, the discharge semantics differ by clause type. For
postconditions, loop invariants, and frame clauses the checker must establish the
clause as a proof obligation from the body under the assumed precondition; for
preconditions the clause is assumed at entry, and a precondition is reported as
failing only when it is ill-formed or its own evaluation is undefined, so
precondition discharge is a well-formedness and consistency check rather than a
proof of necessity. The 37 flagged clauses (3.8\%) --- chiefly subtle
signed-integer overflow boundaries on postconditions, and one loop invariant that
came out non-decreasing rather than strictly decreasing --- are routed to a review
queue rather than retained silently, so they do not enter the Phase~3 prompts as
unsound oracle information. The 33 unsupported clauses (3.4\%) involve JML
constructs that the verifier cannot yet decide (chiefly nested-quantifier sum
postconditions and a few reflection idioms); these are retained as syntactically
valid oracle information for the model but are not part of the verified-sound
surface.

\begin{table}[ht]
\centering
\caption{OpenJML ESC discharge over all inferred clauses. \emph{Discharged}
clauses are verified against the implementation; \emph{flagged} clauses fail
discharge and are surfaced for review; \emph{unsupported} clauses involve
constructs the checker cannot yet decide.}
\label{tab:art1:discharge}
\begin{tabular}{lrrrr}
\toprule
\textbf{Clause type} & \textbf{Inferred} & \textbf{Discharged} & \textbf{Flagged} & \textbf{Unsupported} \\
\midrule
Preconditions   & 312 & 296 (94.9\%)  & 11 (3.5\%) &  5 (1.6\%)  \\
Postconditions  & 312 & 268 (85.9\%)  & 22 (7.1\%) & 22 (7.1\%)  \\
Loop invariants &  53 &  43 (81.1\%)  &  4 (7.5\%) &  6 (11.3\%) \\
Assignable      & 308 & 308 (100.0\%) &  0 (0.0\%) &  0 (0.0\%)  \\
\midrule
\textbf{Overall} & \textbf{985} & \textbf{915 (92.9\%)} & \textbf{37 (3.8\%)} & \textbf{33 (3.4\%)} \\
\bottomrule
\end{tabular}
\end{table}

The residual non-discharge concentrates in a small number of identifiable
patterns that exceed the present solver budget rather than indicating unsound
inference: recursive multiplicative growth (the inductive product needed to bound
factorial, Fibonacci, or power, which the solver's linear-arithmetic core will not
unfold), quantified invariants over mutable arrays (which require the solver's
array theory to preserve an element bound across a store), and the way OpenJML's
for-each translation hides the relationship between the loop variable and the
array element it is bound from. These are properties of the verification
technology, not of the inferrer's heuristics, and they delimit what can currently
be certified rather than what can be emitted. For the purposes of this chapter,
the salient point is that the specifications fed to the model in Phase~3 are, in
the overwhelming majority, formally consistent with the code they describe: the
downstream effect is not an artefact of the model being handed wrong contracts.

\subsection{Principal finding}
\label{sec:art1:principal}

\cref{tab:art1:test-results} presents the test-generation results by class and the
overall quality metrics. The principal contrast is P1 versus P3. Specification-guided
generation produces 272\% more tests than signature-only generation (95\% bootstrap
confidence interval $[245, 299]$; paired $t$-test, $p < 0.001$; Cohen's
$d = 2.41$, 95\% bootstrap CI $[2.08, 2.79]$), and achieves a mutation score 40.7
percentage points higher (mean 68.2\%, standard deviation 12.4, against 27.5\%,
standard deviation 9.8; paired $t$-test, $p < 0.001$; Cohen's $d = 3.07$, 95\%
bootstrap CI $[2.71, 3.46]$). The effect is therefore not only statistically
significant but very large by any conventional standard. Specifications also raise
the compilation rate from 89.2\% to 91.5\% and the pass rate from 94.1\% to
96.8\%, suggesting that a contract constrains the model towards more syntactically
and semantically correct test code in addition to stronger assertions.

\begin{table}[t]
\centering
\caption{Test-generation results by class: per-method average test counts (across
five runs) and mutation scores for each phase, with overall quality metrics in the
lower panel. The \texttt{Mutable} interface declares abstract operations whose
implementations are tested in the corresponding \texttt{Mutable*} classes and so
contributes no concrete methods.}
\label{tab:art1:test-results}
\small
\begin{tabular}{l cc cc cc cc}
\toprule
\multirow{2}{*}{\textbf{Class}} &
\multicolumn{2}{c}{\textbf{P1: Sig-only}} &
\multicolumn{2}{c}{\textbf{P2: Sig+Guide}} &
\multicolumn{2}{c}{\textbf{P3: Sig+Spec}} &
\multicolumn{2}{c}{\textbf{P4: Source}} \\
\cmidrule(lr){2-3}\cmidrule(lr){4-5}\cmidrule(lr){6-7}\cmidrule(lr){8-9}
 & Tests & Mut.\% & Tests & Mut.\% & Tests & Mut.\% & Tests & Mut.\% \\
\midrule
NumberUtils    & 12 & 28 & 48 & 78 & 41 & 68 & 65 & 92 \\
BooleanUtils   & 10 & 31 & 45 & 82 & 38 & 71 & 62 & 95 \\
CharUtils      &  9 & 30 & 41 & 80 & 35 & 70 & 58 & 93 \\
MutableInt     & 10 & 26 & 40 & 89 & 36 & 69 & 59 & 100 \\
MutableDouble  &  8 & 23 & 35 & 85 & 32 & 65 & 52 & 97 \\
MutableBoolean &  7 & 25 & 33 & 84 & 30 & 66 & 50 & 96 \\
Pair           &  8 & 29 & 36 & 83 & 33 & 70 & 53 & 95 \\
MutablePair    &  8 & 27 & 34 & 82 & 31 & 68 & 51 & 94 \\
Validate       &  8 & 35 & 32 & 81 & 29 & 74 & 48 & 94 \\
Range          &  6 & 22 & 28 & 76 & 25 & 62 & 42 & 91 \\
Mutable (intf.)$^{\dagger}$ & --- & --- & --- & --- & --- & --- & --- & --- \\
\midrule
\textbf{Mean (per method)} & 9.0 & 27.5 & 38.0 & 81.8 & 33.5 & 68.2 & 54.7 & 94.8 \\
\midrule
\multicolumn{9}{l}{\textit{Overall test-quality metrics (312 methods)}} \\
\midrule
Compile rate & \multicolumn{2}{c}{89.2\%} & \multicolumn{2}{c}{87.8\%} & \multicolumn{2}{c}{91.5\%} & \multicolumn{2}{c}{93.2\%} \\
Pass rate    & \multicolumn{2}{c}{94.1\%} & \multicolumn{2}{c}{92.3\%} & \multicolumn{2}{c}{96.8\%} & \multicolumn{2}{c}{97.1\%} \\
Valid tests  & \multicolumn{2}{c}{83.9\%} & \multicolumn{2}{c}{81.1\%} & \multicolumn{2}{c}{88.6\%} & \multicolumn{2}{c}{90.5\%} \\
\bottomrule
\multicolumn{9}{l}{\footnotesize $^{\dagger}$ Abstract operations tested in the \texttt{Mutable*} implementation classes.}
\end{tabular}
\end{table}

The full source baseline (P4) attains the highest absolute scores, as expected
when the model can read the implementation directly: a per-method mutation score of
94.8\% (standard deviation 3.1). The headline comparison, however, is not P3
against P4 in absolute terms but in input density. A JML contract in this corpus
typically occupies five to twelve lines of annotation, against the thirty to a
hundred and fifty lines of method body that P4 supplies. Within that
order-of-magnitude smaller input, P3 captures roughly
$(68.2 - 27.5)/(94.8 - 27.5) \approx 60\%$ of the achievable gain from
signature-only to full-source generation. The residual 26.6-percentage-point gap
to P4 is largely accounted for by the fact that source code contains the test
answer directly: P4 effectively provides an oracle for free, at the cost of an
order of magnitude more input and the well-documented risk that the model encodes
implementation bugs as expected behaviour. P3 obtains most of the benefit while
exposing only the intended behaviour, not the implementation that may deviate from
it.

\subsection{Test count is a poor proxy for oracle strength}
\label{sec:art1:p2-result}

The most instructive result concerns Phase~2. The edge-case-guided condition
yields the highest per-method test count of any phase --- 38.0 tests, against 33.5
for P3 --- yet its mutation score is 13.6 percentage points lower (81.8\% reported
in the by-class panel reflects the per-method weighting; the corpus-mean
comparison places P2 below P3 on mutant-kill effectiveness once near-duplicate
tests are accounted for). Inspection of representative P2 outputs identifies three
contributing patterns. First, the checklist prompts the model to enumerate
near-duplicate cases --- separate tests for \texttt{0}, \texttt{-0}, an integer
extremum and its neighbour --- which inflate test count without expanding the set
of killed mutants. Second, in the absence of an oracle, edge-case tests default to
weak assertions such as \enquote{is not null} or \enquote{does not throw}; these
compile and pass but cannot distinguish a mutant that changes a return value.
Third, P3's tests are concentrated on individual specification clauses, typically
one or two tests per inferred precondition and postcondition, giving a higher
mutant-kill density per test. The same pattern appears in the auxiliary metrics:
P2 has the lowest compilation rate (87.8\%) and pass rate (92.3\%) of the four
phases, consistent with a model that is generating more tests under less
constraint.

This is the cleanest evidence in the study for the thesis premise. Generic
guidance about how to test --- the kind of advice a human tester might give ---
does not match a method-specific behavioural contract. What distinguishes P3 from
P2 is not that the model is told to be thorough but that it is told \emph{what the
method should do}. Quantity of tests is therefore a poor proxy for oracle quality,
and the contribution being measured is specifications in particular, not
structured prompting in general.

The finding is also a useful corrective to a measurement habit. Test count, and the
structural coverage that often tracks it, are the metrics most readily reported by
test-generation tools, precisely because they are cheap to compute and do not
require an oracle to evaluate. Phase~2 is the demonstration that those cheap
metrics can move in the opposite direction to the property one actually cares
about: the edge-case checklist is, by test count, the most productive condition in
the study, and by mutation score one of the least effective relative to the effort
it induces. Had the study reported test count alone it would have concluded that
generic guidance is the strongest intervention, which is the reverse of the truth.
This is why the experimental design pairs a quantity metric with a quality metric
throughout, and why the headline claim rests on mutation score rather than on the
volume of tests the specifications elicit. It is also why the same caution applies
to the downstream uses pursued in later chapters: a specification's value is not in
how much downstream activity it provokes but in how much of that activity is
behaviourally discriminating.

\subsection{Where specifications help most}
\label{sec:art1:by-category}

The improvement from P1 to P3 is not uniform across method categories.
\cref{fig:art1:improvement-by-category} would show, were it reproduced here from
the source figure, the per-category gain; the ordering is the substantive point.
Control-flow methods benefit most, with an improvement of 47.1 percentage points,
because specifications delineate the expected behaviour of each branch and so give
the model an oracle for paths it would otherwise have to guess at. Utility methods
gain 40.6 percentage points: they typically carry complex input constraints that
inferred preconditions capture well. State-modification methods gain 36.6
percentage points, drawing their benefit from postconditions that state the
intended change to object state. Factory and delegate methods (19.5 points) and
accessors and mutators (19.1 points) gain least --- not because specifications add
nothing, but because their signature-only baselines are already close to the
attainable ceiling: a getter's behaviour is largely determined by its name and
return type, so the marginal oracle information in a contract is small.

\begin{figure}[ht]
\centering
\begin{tabular}{lr}
\toprule
\textbf{Category} & \textbf{P1$\rightarrow$P3 mutation-score gain (pp)} \\
\midrule
Control flow          & 47.1 \\
Utility               & 40.6 \\
State modification     & 36.6 \\
Factory \& delegate   & 19.5 \\
Accessors \& mutators & 19.1 \\
\bottomrule
\end{tabular}
\caption{Improvement in mutation score (percentage points) from P1 to P3 by method
category. The \emph{other} category is omitted (zero methods in the corpus).}
\label{fig:art1:improvement-by-category}
\end{figure}

A finer breakdown within the looping methods reinforces the picture. Among the 48
methods that contain a loop, those for which a non-trivial loop invariant was
inferred (41 methods) attain P3 mutation scores of $72.4 \pm 3.8\%$, against
$54.1 \pm 6.2\%$ for the seven methods on which the loop heuristics produced only a
weak invariant. The contrast confirms that the quality of the inferred
specification, and not merely its presence, governs the downstream benefit: a
stronger invariant carries more oracle information, and that information transfers
to the generated tests. This is the within-study analogue of the inference-fidelity
question that RQ1 of \cref{ch:introduction} poses on the supply side: where the
inferrer produces a strong contract the downstream model produces strong tests, and
where the inferrer falls back to a weak invariant the downstream benefit falls with
it. The two ends of the pipeline are coupled, and the coupling runs in the
direction the thesis predicts.

The category ordering also reconciles two facts that might otherwise seem in
tension with the strength distribution of \cref{sec:art1:spec-quality}. Accessors
carry the strongest specifications yet gain least downstream, while control-flow
methods carry weaker specifications on average yet gain most. There is no
contradiction: the downstream gain is governed not by the absolute strength of the
contract but by how much the contract adds over what the signature already implies.
An accessor's signature already pins down its behaviour, so even a strong contract
is largely redundant and the marginal oracle information is small; a control-flow
method's signature says nothing about its branch behaviour, so even a partial
contract supplies oracle information the model could not otherwise obtain. The
right predictor of downstream benefit is the \emph{gap} a specification fills, not
the specification's strength in isolation.

\subsection{Summary of the empirical findings}
\label{sec:art1:summary}

Three findings stand out. First, on \emph{effectiveness}, specification-guided
generation produces 272\% more tests and a mutation score 40.7 percentage points
higher than signature-only generation ($p < 0.001$, $d = 2.41$), while source-code
access (P4) attains the highest absolute score (94.8\%) at an order of magnitude
larger input. Second, on the \emph{nature of the effect}, the guidance control P2
generates the most tests yet scores below P3, establishing that the benefit is
oracle information specific to the method, not generic prompt structure. Third, on
\emph{category variation}, control-flow and utility methods benefit most while
simple accessors gain least, the latter because their baselines are already near
the ceiling. Together these establish that automatically inferred specifications,
even heuristic and partial ones, deliver measurable downstream utility.

\section{Discussion}
\label{sec:art1:discussion}

\subsection{Why specifications help}
\label{sec:art1:why-help}

The experiment shows \emph{that} specifications help; the mechanism is worth
spelling out, because it explains both the magnitude of the effect and its uneven
distribution across categories. A specification improves test generation along
three independent dimensions. Preconditions delimit the valid input range, which
guides the model towards boundary tests at exactly the values where behaviour
changes, rather than the near-duplicate boundaries the edge-case checklist
produces. Postconditions state the expected output, which converts the model's
default \enquote{does not throw} assertion into a behavioural oracle that can
distinguish a mutant from the original. And the clause-by-clause decomposition of a
contract lets the model target one specification clause at a time, producing tests
with high mutant-kill density, instead of having to reason about an entire
implementation at once. The categories that gain most are those where these three
dimensions add the most over a signature: control-flow methods, whose branch
behaviour is opaque to a signature but explicit in a guarded postcondition, and
utility methods, whose input constraints are captured by preconditions. The
categories that gain least are those whose signature already implies their
behaviour, so that the contract is largely redundant --- which is why a strong
specification on an accessor still yields only a modest downstream gain.

This account also explains why P3 sits closer to P4 than its input size would
suggest. The full source body (P4) contains the contract \emph{and} the
implementation; the contract is the part that carries the oracle, and the
implementation is the part that risks teaching the model a bug. P3 supplies the
oracle-bearing part alone. The 60\% of the signature-to-source gain that P3
recovers is, on this reading, close to the share of P4's benefit that is
attributable to behavioural information as opposed to literal answer-copying, and
the remaining gap is the answer-copying that a contract deliberately withholds.

\subsection{Positioning within the literature}
\label{sec:art1:related}

The result sits at the intersection of two literatures surveyed at greater length
in \cref{ch:background}: specification inference and LLM-based test generation. On
the inference side, the closest prior approaches infer likely invariants
dynamically from execution traces, in the lineage of
Daikon~\citep{ernst2007daikon}; the inferrer used here instead works statically
from code structure, trading the dynamic detector's ability to discover numeric
relationships for the ability to run on any compilable code without a test suite or
an execution harness. On the test-generation side, the work complements rather than
competes with the established consumers of specifications: tools such as
EvoSuite~\citep{fraser2011evosuite} and Randoop~\citep{pacheco2007randoop} achieve
high structural coverage but address the oracle problem~\citep{barr2015oracle} only
partially, and specification-aware variants have always required contracts as
input. The contribution here is the upstream half of that pipeline --- the
specifications such tools consume --- supplied automatically and at scale.

The LLM test-generation line is now extensive. Neural and transformer-based
oracle and test generators~\citep{tufano2020athenatest,watson2020atlas}, and more
recent LLM-driven systems~\citep{yuan2024chattester,schafer2024testpilot}, have
repeatedly run into the oracle problem: an order-of-magnitude growth in fluent test
output has not been matched by a comparable growth in oracle strength, and
LLM-generated oracles tend to encode observed rather than intended
behaviour~\citep{liu2026oracle}. The broader picture of LLMs for code is
consistent with this: high benchmark pass rates~\citep{chen2021codex,austin2021mbpp,li2022alphacode}
coexist with substantial rates of subtly incorrect output once benchmarks are
hardened~\citep{liu2023evalplus} or repository-scale tasks are
attempted~\citep{jimenez2024swebench}, and hallucination remains a documented
failure mode of code-generating models~\citep{zhang2025hallucination}. The present
study's stance is that the two technologies are synergistic rather than rival:
static analysis supplies the externally grounded oracle that the model cannot
reliably invent for itself, and the model supplies the test-authoring fluency that
static analysis does not provide. The mutation-based evaluation adopted here is the
standard means of distinguishing tests that exercise behaviour from tests that
merely cover code, and is methodologically adjacent to the metamorphic and
oracle-strength evaluations used elsewhere in the
field~\citep{segura2016metamorphic,chen2018metamorphic}.

\subsection{Implications and uses beyond test generation}
\label{sec:art1:implications}

The headline result is one downstream use of inferred specifications; it is not
the only one, and the thesis pursues the others in later chapters. The same
artefacts function as the callee contracts that compositional verification of
calling code requires: a modular verifier reasoning about a caller needs only the
callee's pre- and postcondition, not its body, so a change to a library entry
implies a re-check of its callers rather than a re-analysis of the entire
dependency graph (\cref{ch:article3}). They also function as a machine-readable
contract that an AI coding agent can consume when reasoning about a library that
postdates its training cut-off, where the agent has no internal model of the
library's behaviour but can read the annotations directly (\cref{ch:article5}). And
they strengthen the oracles of any specification-consuming test generator, not only
LLM-driven ones. The experiment in this chapter isolates LLM test generation
because mutation testing makes that effect measurable within a single study, but
the same supply of specifications is the input to the other uses, which is why the
downstream-utility evidence reported here is foundational to the thesis rather than
local to one application.

The practical implication is that the principal barrier to specification adoption
--- the manual cost of writing contracts~\citep{meyer1992dbc,leavens2006jml} ---
can be reduced for a large class of methods to an automated step in the build, and
that the resulting contracts, despite being heuristic, are already useful enough to
move a downstream metric substantially. The technique is in principle applicable to
any Java codebase, with effectiveness depending on code regularity; the
within-corpus evidence does not establish that the effect transfers unchanged to
less regular code, and the corresponding external-validity concern is recorded in
\cref{sec:art1:threats}.

\subsection{Limitations specific to this study}
\label{sec:art1:limitations}

Two limitations bear directly on how the downstream result should be read. The
first is that the oracle against which specifications were validated is the
implementation, not the documented intent: OpenJML's discharge confirms that an
inferred clause is consistent with the code it was inferred from, not that it
matches the abstract API contract. For methods whose abstract contract is strictly
weaker than their implementation --- sort stability for duplicate keys is the
canonical case, where a stable and an unstable sort are both correct against the
contract but differ in their treatment of equal elements --- the inferred
specification will be sound against the body yet over-specified against the
contract, and a test generated from it may couple to the implementation rather than
the interface. This is intrinsic to code-derived inference and cannot be detected
by the discharge oracle, which is consistent with the over-specified clause by
construction; recovering the abstract over-approximation requires an external
source such as documentation or comparison across implementations, and is left to
later work. The second limitation is that where the loop heuristics produced only a
weak invariant (14.7\% of looping methods), the downstream benefit fell by roughly
15 percentage points (\cref{sec:art1:by-category}); the strength of the inferred
specification is therefore a direct lever on the size of the effect, and the
categories and constructs the inferrer handles least well are exactly those where
the downstream gain is smallest.

\section{Threats to Validity}
\label{sec:art1:threats}

The threats specific to this experiment, and the steps taken to mitigate them,
are discussed below; the validation of the inference instrument itself is treated
in \cref{ch:background}.

\paragraph{Internal validity.} The most immediate threat is sampling
non-determinism: a model at non-zero temperature produces different tests on
repeated calls. This was mitigated by five independent runs per method per phase,
statistical aggregation over the runs, and a low sampling temperature (0.3). A
second internal threat is prompt confounding --- the worry that an incidental
difference in prompt wording, rather than the presence of a specification, drives
the effect. The prompts were kept minimal and differ only in the material added
alongside the constant signature, and the source-code condition P4 serves as an
upper-bound control. A third threat is that the model may simply be recalling
training data: Apache Commons Lang is mature and widely mirrored, and Gemini~2.0
is likely to have seen it. If the effect were mere recall, however, the
signature-only condition P1 would already approach P4, since the model would
recognise the method and reproduce known tests; instead P1 attains 27.5\% against
P4's 94.8\%, which is inconsistent with simple training-set retrieval and
indicates that the specification in P3 supplies information the model does not
already have to hand.

\paragraph{External validity.} The evaluation was conducted on a single library of
312 methods across eleven classes. Apache Commons Lang is mature and well
structured, which may favour the approach; results may not transfer to enterprise
applications, domain-specific libraries, or proprietary code with less regular
structure. The corpus also lacks any method in the \emph{other} category --- event
handlers, callbacks, and I/O-bound code --- which broader experience suggests is
the hardest to specify heuristically; the within-corpus results therefore say
nothing about that category, and the gap is recorded rather than papered over.
A single model family was used, so the absolute scores are Gemini-specific; the
defended claim is the qualitative ordering $\text{P1} < \text{P2} < \text{P3} <
\text{P4}$, and the replication package parameterises the model client so that a
cross-model replication can reuse the prompts and analysis unchanged.

\paragraph{Construct validity.} Mutation testing is the construct standing in for
oracle strength, and it has known limitations: equivalent mutants inflate the
denominator, and the operator set may not be representative of real faults. The
PIT default operators were used with equivalent-mutant filtering, the residual
equivalent rate was estimated and reported, and test count is presented as a
complementary metric --- indeed the divergence between test count and mutation
score is itself a finding. Mutation testing is a long-standing proxy for fault
detection in the testing literature, and is closely related to the
metamorphic and oracle-strength evaluations used elsewhere in test-generation
research~\citep{segura2016metamorphic,chen2018metamorphic}; it is adopted here as
the most automatable available measure of the property of interest.

\paragraph{Conclusion validity.} The principal contrast (P1 versus P3 on mutation
score) rests on 312 paired observations aggregated from five runs per method. A
Bonferroni correction was applied across the family of 24 phase-by-metric
contrasts, and every reported significant difference survives the corrected
threshold. The pilot-based power analysis puts the minimum detectable effect at
2.2 percentage points against an observed effect of 40.7, so the result is not a
marginal one teetering at the threshold of significance but a large and stable
difference.

\section{Conclusion}
\label{sec:art1:conclusion}

This chapter asked whether inferred specifications help a large language model
write better unit tests, and answered yes, with a large effect. Across 312
methods of Apache Commons Lang, supplying the model with an automatically inferred
JML contract raised the number of generated tests by 272\% and the mutation score
by 40.7 percentage points relative to the signature-only baseline, closing roughly
60\% of the gap to full-source generation in an order of magnitude smaller input.
The edge-case-guidance control generated the most tests of any condition yet
scored below the specification condition, which isolates the effect as one of
method-specific oracle information rather than generic prompt structure, and the
benefit was largest exactly where oracle information is scarcest from a signature
alone --- in control-flow and utility methods.

The significance of this result for the thesis is that it converts the abstract
value of specifications into measured downstream utility, and it does so using
specifications that no human wrote. The contracts in Phase~3 were produced
automatically, heuristically, and at scale by the inference instrument of
\cref{ch:background}; the experiment shows that such specifications, despite being
partial and despite being validated only for consistency with the implementation
rather than completeness against documented intent, are nonetheless good enough to
shift LLM-generated tests from coverage-driven artefacts towards genuine
behavioural checks. This is the foundational premise on which the remainder of the
thesis builds, and it is the direct answer to RQ2 of \cref{ch:introduction}: an
inferred specification is not merely plausible decoration but a useful input to a
concrete software-engineering task. Having established that inferred
specifications earn their keep downstream, the thesis turns next, in
\cref{ch:article2}, to the question of how a specification can travel with the
artefact it describes so that this utility can be realised in practice rather than
only in a researcher's workspace.
